\section{Adaptation to all four budgets}
\label{app:stochastic-adaptation}

This appendix removes the stochastic-budget inputs from
Theorem~\ref{thm:approximation-adaptation}. An independent block of
covariates supplies two empirical local Rademacher fixed points. The
construction follows the empirical fixed-point method of
\citet{bartlett2005local}. Two additional steps are needed for the present
experiment: controlling the star hull of an arbitrary represented
difference core without multiplying its declared radius by a logarithm,
and transferring a resulting linear core certificate to the full outer
complexity under sequential M1. Neither step assumes that a public class
is star-shaped.

\begin{theorem}[Adaptation to all four budgets]
\label{thm:full-budget-adaptation}
Fix a represented public instance $\Gamma$, the regularity tuple $\rho$,
and $n\ge1$. There is one Borel rule
$\delta_{\Gamma,\rho,n}:\mathcal Z^{2n}\to[-3,3]$ such that,
simultaneously for every $(a,s,b,t)\in[0,1]^4$ and
$P\in\mathcal P_{\rho,n,a,s,b,t}(\Gamma)$,
\begin{equation}
 \begin{split}
 \E_P|\delta_{\Gamma,\rho,n}-\beta(P)|
 &\le C_\rho\left[q+
 \min\{a(b+t)+a^2+s^2,\ b(a+s)+(b+t)^2\}\right]\\
 &\le C_\rho(q+W).
 \end{split}
 \label{eq:full-budget-adaptation-rate}
\end{equation}
The rule uses $\Gamma,\rho,n$ and the observations, but none of
$a,s,b,t$, the covariate law, or a population comparator. The constant
is independent of the public instance, the budgets, the law, and $n$.
Only the average lower bound $\E u^2\ge v_-$ is required.
\end{theorem}

\subsection{A star-hull bound for an arbitrary countable class}
\label{app:stochastic-star}

For a countable symmetric class $\mathcal H$ containing zero, write
\[
 \operatorname{star}(\mathcal H)
 =\{\alpha f:f\in\mathcal H,\ 0\le\alpha\le1\}.
\]
Every localized supremum over this class equals the corresponding
supremum with $\alpha\in\mathbb Q\cap[0,1]$. Indeed, any permitted
real coefficient can be approximated from below by rational
coefficients, preserving either a population or an empirical upper
norm constraint and converging in every finite-sample linear process.
Squares and population quadratic means are also continuous in the
scaling coefficient, so the same argument makes the square-process
suprema and ball-inclusion events below measurable. We may use the real star hull
for its scaling property and the countable rational star hull for
measurability, without changing a supremum.

\begin{lemma}[Cost of the star hull above the logarithmic floor]
\label{lem:stochastic-star-cost}
Let $\mathcal H$ be a countable symmetric class containing zero, with
envelope one. Let $m\ge1$, $h_m=\sqrt{\log(em)/m}$, and $s\ge0$.
Suppose that, for a fixed $\kappa\ge0$,
\[
 \mathfrak R_{m,P_X}(r;\mathcal H)\le\kappa sr,
 \qquad r\ge s,\quad r>0.
\]
There is a constant $C_\star$, depending only on $\kappa$, such that
\begin{equation}
 \mathfrak R_{m,P_X}(r;\operatorname{star}(\mathcal H))
 \le C_\star(s+h_m)r,
 \qquad r\ge s\vee h_m.
 \label{eq:stochastic-star-cost}
\end{equation}
No closure of $\mathcal H$ under scaling is assumed.
\end{lemma}
\begin{proof}
Fix $s\vee h_m\le r<1$ and put
\[
 J=\lceil\log_2(1/r)\rceil,\qquad R_j=2^jr,
 \qquad
 Z_j=\sup_{\substack{f\in\mathcal H\\\|f\|_2\le R_j}}
 \left|\frac1m\sum_{i=1}^m\epsilon_i f(X_i)\right|,
 \quad 0\le j\le J.
\]
If $\alpha f$ has population norm at most $r$, a function with
$\|f\|_2\le r$ contributes at most $Z_0$. A function with
$2^{j-1}r<\|f\|_2\le2^jr$ has $\alpha\le2^{1-j}$.
Consequently, pointwise in the sample and signs, its star-hull
supremum is at most
\begin{equation}
 \max\{Z_0,\ \max_{1\le j\le J}2^{1-j}Z_j\}.
 \label{eq:stochastic-star-max}
\end{equation}

Apply Theorem 2.1 of \citet{bartlett2005local} to the signed
observations $(X_i,\epsilon_i)$ and the centered symmetric class
$\{(x,e)\mapsto ef(x):\|f\|_2\le R_j\}$. Its range is
$[-1,1]$, its variance is at most $R_j^2$, and its Rademacher mean
is exactly $\mathfrak R_{m,P_X}(R_j;\mathcal H)$: multiplying
$\epsilon_i$ by a second independent Rademacher sign preserves the
joint sign distribution. With the source theorem's tuning parameter
equal to one, it follows that
\[
 \Pr\left\{Z_j>4\kappa sR_j+
 R_j\sqrt{\frac{2x_j}{m}}+\frac{8x_j}{3m}\right\}
 \le e^{-x_j}.
\]
Take $x_j=y+2j+\log4$, for $y\ge0$, and sum these failures.
After multiplying the outer-shell bounds by $2^{1-j}$, the maximum
in \eqref{eq:stochastic-star-max} is at most
\[
 C\kappa sr+
 Cr\sqrt{\frac{y+J+1}{m}}+\frac{C(y+1)}m
\]
except with probability at most $e^{-y}$. Here
$\sup_{j\ge1}2^{-j}(y+2j+\log4)\le C(y+1)$.
Integrating this tail gives
\[
 \mathfrak R_{m,P_X}(r;\operatorname{star}(\mathcal H))
 \le C\kappa sr+Cr\sqrt{\frac{J+1}{m}}+\frac Cm.
\]
Since $r\ge h_m$, we have $J+1\le C\log(em)$ and
$m^{-1}\le rh_m$. This proves
\eqref{eq:stochastic-star-cost} for $r<1$. For $r\ge1$ the
localized core and star-hull suprema are their global suprema, which
are equal; the retained condition $r\ge s\vee h_m$ allows the
original certificate to prove the assertion directly.
The argument includes $s=0$. It takes the expectation of a maximum
using concentration; summing the means of the shells would give an
unnecessary logarithmic factor.
\end{proof}

\subsection{An empirical radius and its population certificate}
\label{app:stochastic-calibration}

For an iid calibration sample $X_1,\ldots,X_m$, let $Q_m$ denote
its empirical average. Auxiliary Rademacher signs below are averaged
exactly over their $2^m$ possible values; no additional randomization
is part of the estimating rule.

\begin{lemma}[Measurable calibration of a countable difference core]
\label{lem:stochastic-empirical-radius}
Let $\mathcal F$ be a countable symmetric class containing zero,
with envelope at most $L$, where $L\ge1$ is known. Suppose that
\begin{equation}
 \mathfrak R_{m,P_X}(u;\mathcal F)
 \le\kappa su,\qquad u\ge s,\quad u>0,
 \label{eq:stochastic-core-input}
\end{equation}
for some $s\ge0$ and fixed $\kappa$. Set
$\overline{\mathcal F}=\{f/L:f\in\mathcal F\}$,
$x=6\log(em)$, and, for $r>0$, define
\begin{align}
 \widehat\Phi(r)
 &=2^{-m}\sum_{e\in\{-1,1\}^m}
 \sup_{f\in\overline{\mathcal F}}
 \left(1\wedge\sqrt{\frac{2r}{Q_mf^2}}\right)
 \left|\frac1m\sum_{i=1}^m e_if(X_i)\right|,
 \label{eq:stochastic-empirical-phi}\\
 \widehat\psi(r)&=20\widehat\Phi(r)+\frac{31x}{m}.
 \label{eq:stochastic-empirical-psi}
\end{align}
A summand with $Q_mf^2=0$ is defined to be zero. Let
\begin{equation}
 \widehat r
 =\inf\{r\in\mathbb Q:r>0,\ \widehat\psi(r)\le r\},
 \qquad \widehat R=L\sqrt{\widehat r}.
 \label{eq:stochastic-radius-rule}
\end{equation}
This defines a finite positive Borel statistic, using neither $s$ nor
$P_X$. With probability at least $1-4e^{-x}$, simultaneously,
\begin{align}
 \mathfrak R_{m,P_X}(u;\mathcal F)
 &\le\frac{\widehat R}{10}\,u
 &&\text{for every }u\ge\widehat R,
 \label{eq:stochastic-radius-certificate}\\
 \widehat R&\le C_{L,\kappa}(s+h_m).
 \label{eq:stochastic-radius-size}
\end{align}
Moreover $\widehat r\ge31x/m$ on every sample.
\end{lemma}
\begin{proof}
Write $\mathcal S=\operatorname{star}(\overline{\mathcal F})$.
The expression in \eqref{eq:stochastic-empirical-phi} is exactly
\[
 \widehat\Phi(r)
 =\E_\epsilon\sup_{\substack{f\in\mathcal S\\Q_mf^2\le2r}}
 \left|\frac1m\sum_i\epsilon_if(X_i)\right|.
\]
It is Borel for each $r>0$, by its finite sign average and countable
supremum. Each term before taking the supremum, as a function of
$\sqrt r$, is Lipschitz with constant at most $\sqrt2$: when
$Q_mf^2>0$, divide its linear process by $\sqrt{Q_mf^2}$ and use
Cauchy--Schwarz. Thus $\widehat\Phi$ is continuous on $(0,\infty)$.
It is nondecreasing and $\widehat\Phi(r)/\sqrt r$ is nonincreasing.
The same properties hold for $\widehat\psi$, which is positive
because of its deterministic additive term. It follows that
$\widehat\psi$ has a unique positive fixed point and that
$\widehat\psi(r)\le r$ exactly when $r$ is at least this fixed
point. In particular \eqref{eq:stochastic-radius-rule} equals the
fixed point and is Borel. The bounds
\[
 \frac{31x}{m}\le\widehat r\le20+\frac{31x}{m}
\]
follow from $0\le\widehat\Phi(r)\le1$. These observations also
handle functions with zero empirical norm.

For the lower comparison, define the nonrandom function
\[
 \Phi(r)=\mathfrak R_{m,P_X}(\sqrt r;\mathcal S),\qquad
 \psi_1(r)=10\Phi(r)+\frac{11x}{m},
\]
and let $r_1$ be the unique positive fixed point of $\psi_1$.
Scaling in $\mathcal S$ shows directly that $\Phi$ is nondecreasing
and $\Phi(r)/\sqrt r$ is nonincreasing. These properties also give
continuity on $(0,\infty)$ and the asserted fixed-point properties
of $\psi_1$. We use three primitive results of
\citet{bartlett2005local}: Corollary 2.2, Lemma 3.6, and Lemma A.4,
always with envelope one. Their suprema are measurable by the
rational-star observation above.

Corollary 2.2 at the fixed, population-dependent radius $r_1$ gives
\[
 \{f\in\mathcal S:Pf^2\le r_1\}
 \subseteq\{f\in\mathcal S:Q_mf^2\le2r_1\}
\]
except with probability $e^{-x}$, because
$r_1=10\Phi(r_1)+11x/m$. Apply the lower comparison in Lemma A.4
with its tuning parameter $1/2$ to the fixed class on the left.
Its range is $[-1,1]$, so, except with probability $e^{-x}$,
\[
 \Phi(r_1)
 \le2\E_\epsilon
 \sup_{\substack{f\in\mathcal S\\Pf^2\le r_1}}
 \left|\frac1m\sum_i\epsilon_if(X_i)\right|+
 \frac{2x}{m}.
\]
Symmetry makes the absolute supremum equal to the source's
one-sided supremum. On these two events,
\[
 r_1=10\Phi(r_1)+\frac{11x}{m}
 \le20\widehat\Phi(r_1)+\frac{31x}{m}
 =\widehat\psi(r_1),
\]
and therefore $r_1\le\widehat r$. For every
$v\ge\sqrt{\widehat r}$, the sub-root property implies
\[
 \mathfrak R_{m,P_X}(v;\mathcal S)
 =\Phi(v^2)\le\frac{\psi_1(v^2)}{10}
 \le\frac{\sqrt{r_1}}{10}v
 \le\frac{\sqrt{\widehat r}}{10}v.
\]
Since $\overline{\mathcal F}\subseteq\mathcal S$, rescaling by
$L$ yields, for $u\ge L\sqrt{\widehat r}$,
\[
 \mathfrak R_{m,P_X}(u;\mathcal F)
 \le\frac{\sqrt{\widehat r}}{10}u
 \le\frac{\widehat R}{10}u,
\]
which proves \eqref{eq:stochastic-radius-certificate}.

For the upper comparison, \eqref{eq:stochastic-core-input} and
$L\ge1$ imply
\[
 \mathfrak R_{m,P_X}(v;\overline{\mathcal F})
 \le\kappa sv,\qquad v\ge s,\quad v>0.
\]
Lemma~\ref{lem:stochastic-star-cost} supplies a constant $C_\star$
such that, with $w=s+h_m$,
\[
 \Phi(r)\le C_\star w\sqrt r
 \qquad\text{whenever }\sqrt r\ge s\vee h_m.
\]
Freeze
\begin{equation}
 D=[120(C_\star+1)]^2,\qquad r_0=Dw^2.
 \label{eq:stochastic-proof-radius}
\end{equation}
These quantities are proof devices; the statistic uses neither of
them. Since $x/m=6h_m^2\le6w^2$, this choice ensures
\[
 2r_0\ge20\Phi(2r_0)+\frac{26x}{m}.
\]
Lemma 3.6 at $2r_0$ consequently gives
\[
 \{f\in\mathcal S:Q_mf^2\le2r_0\}
 \subseteq\{f\in\mathcal S:Pf^2\le4r_0\}
\]
except with probability $e^{-x}$. The upper comparison in Lemma A.4,
again with tuning parameter $1/2$ and range $[-1,1]$, applied to
the fixed population ball on the right gives, except with
probability $e^{-x}$,
\[
 \E_\epsilon
 \sup_{\substack{f\in\mathcal S\\Pf^2\le4r_0}}
 \left|\frac1m\sum_i\epsilon_if(X_i)\right|
 \le\frac32\Phi(4r_0)+\frac{4x}{3m}.
\]
On these two events,
\[
 \begin{split}
 \widehat\psi(r_0)
 &\le30\Phi(4r_0)+\frac{173x}{3m}\\
 &\le\{60C_\star\sqrt D+346\}w^2
 \le Dw^2=r_0.
 \end{split}
\]
Thus $\widehat r\le r_0$ and
$\widehat R\le L\sqrt D(s+h_m)$. A union bound over the four
events proves the lemma. This argument uses only the stated ball
inclusions and Rademacher comparisons. In particular, it does not
require a loss-class assumption $Pf^2\le BPf$ or a lower bound on
an unregularized population fixed point. The positive confidence
term handles that fixed point being arbitrarily small or zero.
\end{proof}

\subsection{Sequential M1 transfers a linear certificate to the full class}
\label{app:stochastic-full-transfer}

\begin{lemma}[Transfer with a positive radius margin]
\label{lem:stochastic-full-transfer}
For a public class satisfying sequential M1 and the common envelope,
every sample size $m\ge1$, every $u\ge0$, and every $\eta>0$ satisfy
\begin{equation}
 \mathfrak R^*_{m,P_X}(u;\Delta\mathcal G)
 \le\mathfrak R_{m,P_X}(u+\eta;\Delta\mathcal G^{\rm rep}).
 \label{eq:stochastic-margin-transfer}
\end{equation}
Consequently, if for some $R\ge0$ and $c\ge0$ the core has the
linear certificate
\[
 \mathfrak R_{m,P_X}(v;\Delta\mathcal G^{\rm rep})\le cRv,
 \qquad v\ge R,\quad v>0,
\]
then the same certificate holds for the full outer complexity at
every $u\ge R$ with $u>0$.
\end{lemma}
\begin{proof}
Fix $f=g-g'\in\Delta\mathcal G$ with $\|f\|_2\le u$.
Sequential M1 gives represented sequences converging pointwise
everywhere to $g$ and $g'$. Their differences $f_k$ converge
pointwise everywhere to $f$ and, by the common envelope and
dominated convergence, in $L_2(P_X)$. Thus
$\|f_k\|_2\le u+\eta$ eventually. For every covariate sample
and every sign vector, the signed sample average of $f_k$ converges
to that of $f$. Its limit in absolute value is bounded by the core
supremum at radius $u+\eta$. Taking the supremum over all such
full-class $f$ proves the pointwise inequality
\[
 Z_{\Delta\mathcal G,m}(u)
 \le Z_{\Delta\mathcal G^{\rm rep},m}(u+\eta).
\]
The right side is measurable and is therefore an admissible
majorant in the definition of outer expectation. This proves
\eqref{eq:stochastic-margin-transfer}. Under the linear certificate,
the right side in that display is at most $cR(u+\eta)$; let
$\eta\downarrow0$. No equality of full and represented localized
complexities at the same radius is asserted.
\end{proof}

\subsection{The three-block rule and its risk}
\label{app:stochastic-wrapper}

We now prove Theorem~\ref{thm:full-budget-adaptation}. All constants
of Appendix~\ref{app:approximation-adaptation} are frozen for the
known regularity tuple
\[
 \rho'=(B_0,B_{\mathcal G},v_-,v_+,2c_{\rm rad}).
\]
Let $m_{\rho'}$ be its fixed lower sample-size threshold, enlarged
to at least two. If $m=\lfloor2n/3\rfloor<m_{\rho'}$, define the
rule to be zero. There are only finitely many such $n$, with a
bound depending only on $\rho$, so its absolute loss, at most
three, is absorbed by $C_\rho q$. This includes $n=1$.

For the remaining $n$, split the first $3m\le2n$ observations
deterministically into three disjoint blocks, each of size $m$.
The first is the calibration block; its treatments and outcomes
are not used. The second and third supply the two input halves of
the rule of Theorem~\ref{thm:approximation-adaptation} at sample
size $m$. Any leftover observations are unused. In particular,
$m\le n$ and $n/m\le2$.

For each $\nu\in\{\mu,\pi\}$, apply
\eqref{eq:stochastic-empirical-phi}--\eqref{eq:stochastic-radius-rule}
on the calibration covariates to
$\mathcal F_\nu=\Delta\mathcal G_\nu^{\rm rep}$, with
$L=\max\{1,2B_{\mathcal G}\}$. Denote its output by
$\widehat R_\nu$. Define
\begin{equation}
 C_{\rm cap}=\max\{1,(60c_{\rm rad})^{-1}\},\qquad
 S=\min\{1,C_{\rm cap}\widehat R_\mu\},\qquad
 T=\min\{1,C_{\rm cap}\widehat R_\pi\}.
 \label{eq:stochastic-capped-inputs}
\end{equation}
The final rule is the explicit estimator of
Appendix~\ref{app:approximation-adaptation}, with public instance
$\Gamma$, regularity tuple $\rho'$, sample size $m$, and stochastic
inputs $(S,T)$, applied to the last two blocks. Here $T$ in
\eqref{eq:stochastic-capped-inputs} denotes the scalar stochastic
input, not a treatment observation.

The estimator just specified is jointly Borel in its two stochastic
inputs and its observations. To see this directly from Appendix~
\ref{app:approximation-adaptation}, its initial fits and pilot are
Borel and do not use those inputs; each calibration tolerance is a
continuous function of their squares; and its constraints, countable
feasible-set decisions, countable infima, and first strict approximate
minimizers are Borel jointly in inputs and observations. The constants
are fixed for $\rho'$ before either input is selected. Composition
with \eqref{eq:stochastic-radius-rule} and
\eqref{eq:stochastic-capped-inputs} therefore gives one Borel rule
on $\mathcal Z^{2n}$. Its definition contains none of the four
declared budgets.

Fix now any $(a,s,b,t)$ and compatible $P$, solely for the proof.
For every fixed population-localized countable core, adding $n-m$
independent signed observations and applying conditional Jensen
gives
\begin{equation}
 \mathfrak R_{m,P_X}(u;\mathcal F_\nu)
 \le\frac nm\mathfrak R_{n,P_X}(u;\mathcal F_\nu).
 \label{eq:stochastic-sample-transfer}
\end{equation}
Indeed, condition on the first $m$ covariates and signs; the added
signed summands have conditional mean zero, and the supremum of
absolute linear averages is convex. Countability permits ordinary
conditional expectation. Combining this with core inclusion into
the full class and the original certificate gives
\[
 \mathfrak R_{m,P_X}(u;\mathcal F_\mu)\le6c_{\rm rad}su
 \quad(u\ge s,\ u>0),
\]
and the corresponding inequality with $t$ for $\mathcal F_\pi$.
Thus Lemma~\ref{lem:stochastic-empirical-radius} applies with
$\kappa=6c_{\rm rad}$. There is an event $E$ measurable with
respect to the calibration covariates, of probability at least
$1-8(em)^{-6}$, on which both output certificates hold and
\begin{equation}
 \widehat R_\mu\le C_\rho(s+h_m),\qquad
 \widehat R_\pi\le C_\rho(t+h_m).
 \label{eq:stochastic-both-sizes}
\end{equation}
One may take $E$ to be the intersection of the four fixed-radius
events constructed for each core in the proof of that lemma. This
description makes its measurability explicit, without a supremum
over random population radii.

On $E$, Lemma~\ref{lem:stochastic-full-transfer} turns each output
core certificate into the same full outer certificate. If $S<1$,
then $S=C_{\rm cap}\widehat R_\mu\ge\widehat R_\mu$, so for
every $u\ge S$,
\[
 \mathfrak R^*_{m,P_X}(u;\Delta\mathcal G_\mu)
 \le\frac{\widehat R_\mu}{10}u
 \le6c_{\rm rad}Su=3c'_{\rm rad}Su.
\]
If $S=1$, use $s\le1$ in the original certificate and
\eqref{eq:stochastic-sample-transfer}: for every $v\ge1$,
\[
 \mathfrak R_{m,P_X}(v;\mathcal F_\mu)
 \le6c_{\rm rad}sv\le6c_{\rm rad}v.
\]
The positive-margin transfer of
Lemma~\ref{lem:stochastic-full-transfer}, followed by letting the
margin tend to zero, gives the required full outer bound
$6c_{\rm rad}u$ for every $u\ge1$. The same two cases apply to
$T$ and the propensity class. This argument uses ordinary Jensen
only for countable cores; it does not require a conditional-Jensen
identity for nonmeasurable full-class suprema.

It follows that, for every calibration realization in $E$,
\begin{equation}
 P\in\mathcal P_{\rho',m,a,S,b,T}(\Gamma).
 \label{eq:stochastic-conditional-certificate}
\end{equation}
The approximation budgets in this assertion are the original $a,b$:
neither public class, either approximation infimum, nor the law has
changed. Only the sample size, the stochastic certificates, and the
fixed complexity constant have changed. Also, from
\eqref{eq:stochastic-both-sizes},
\[
 S\le C_\rho(s+h_m),\qquad T\le C_\rho(t+h_m)
 \quad\text{on }E.
\]

Condition on the calibration covariates. Independence of the three
blocks leaves the other two blocks iid with joint law $P^{2m}$.
For realizations in $E$, apply
Theorem~\ref{thm:approximation-adaptation} using
\eqref{eq:stochastic-conditional-certificate}. Its conditional
expected absolute loss is bounded by
\[
 C_\rho\left[m^{-1/2}+
 \min\{a(b+T)+a^2+S^2,\ b(a+S)+(b+T)^2\}\right].
\]
Quadratic absorption gives, on this same event,
\begin{align*}
 a(b+T)+a^2+S^2
 &\le C_\rho\{a(b+t)+a^2+s^2+h_m^2\},\\
 b(a+S)+(b+T)^2
 &\le C_\rho\{b(a+s)+(b+t)^2+h_m^2\}.
\end{align*}
For the first inequality use $ah_m\le(a^2+h_m^2)/2$.
For the second use $bh_m\le(b^2+h_m^2)/2$ and
$b^2\le(b+t)^2$. Since $m$ is a fixed proportion of $n$,
\[
 m^{-1/2}\le Cq,\qquad
 h_m^2=\frac{\log(em)}m\le Cq.
\]
Both the output and the target lie in $[-3,3]$, so the contribution
from $E^c$ is at most $48(em)^{-6}\le Cq$. Averaging proves the
first inequality in \eqref{eq:full-budget-adaptation-rate}.
The second follows from $a^2+s^2\le(a+s)^2$ and the definition
of $W$.

All radii in the calibration rule are strictly positive because of
the confidence term. Zero declared $s$ or $t$ enters the proof only
through its original positive-radius certificate, which then
vanishes; Lemma~\ref{lem:stochastic-star-cost} and all subsequent
steps still apply. No radius-zero quotient, additional variance
assumption, or fitted-index substitution into an unproved process
event occurs. This completes the proof of
Theorem~\ref{thm:full-budget-adaptation}.
